\section{Learning to learn with recurrent neural networks}
\label{sec:learning-to-learn}

% Consider the minimization operator $\calM=\min_\theta$ which can be applied to
% a function $f$ and returns its minimum value. The standard problem in
% optimization is then that of solving $\calM(f)$. However this operator is in
% some sense an idealized oracle. Instead, in this work we will parameterize this
% operator $\calM_\phi$ and ask, for a given distribution of functions $f$ what
% is the quality of this operator, or

In this work we consider directly parameterizing the optimizer. As a result, in
a slight abuse of notation we will write the final \emph{optimizee parameters}
$\theta^*(f, \phi)$ as a function of the \emph{optimizer parameters} $\phi$ and
the function in question. We can then ask the question: What does it mean for
an optimizer to be good? Given a distribution of functions $f$ we will write
the expected loss as
\begin{equation}
    \calL(\phi) =
    \mathbb{E}_{f} \Big[ f\big(\theta^*(f, \phi)\big) \Big] \,.
    \label{eq:meta-optimization-final}
\end{equation}
As noted earlier, we will take the update steps $g_t$ to be the output of a
recurrent neural network $m$, parameterized by $\phi$, whose
state we will denote explicitly with $h_t$. Next, while the objective function
in \eqref{eq:meta-optimization-final} depends only on the final parameter
value, for training the optimizer it will be convenient to have an objective
that depends on the entire trajectory of optimization, for some horizon T,
\begin{equation}
    \calL(\phi)
    = \mathbb{E}_{f} \left[\sum_{t=1}^T w_t f(\theta_t)\right]
    \qquad\text{where}\qquad
    \begin{aligned}[t]
        \theta_{t+1} &= \theta_t + g_t \,,
        \\
        \begin{bmatrix}
            g_t \\ h_{t+1}
        \end{bmatrix}
        &= m(\nabla_t, h_t, \phi) \,.
    \end{aligned}
    \label{eq:meta-optimization}
\end{equation}
Here $w_t \in \mathbb{R}_{\ge 0}$ are arbitrary weights associated with each
time-step and we will also use the notation $\nabla_t=\nabla_\theta
f(\theta_t)$. This formulation is equivalent to
\eqref{eq:meta-optimization-final} when $w_t = \boldsymbol{1}[t = T]$,
but later we will describe why using different weights can prove useful.
%Later it will be useful to relax this weighting in order to train the optimizer
%efficiently, but first we will discuss how to train the optimizer when the
%$w_t$ are arbitrary non-negative weights on the trajectory of iterates.

We can minimize the value of $\mathcal{L}(\phi)$ using gradient descent on
$\phi$. 
The gradient estimate
$\partial\mathcal{L(\phi)}/\partial\phi$ can be computed by sampling a random
function $f$ and applying backpropagation to the computational graph in
Figure~\ref{fig:graph}.  We allow gradients to flow along the solid edges in
the graph, but gradients along the dashed edges are dropped.  Ignoring
gradients along the dashed edges amounts to making the assumption that the
gradients of the optimizee do not depend on the optimizer parameters, i.e.\
$\partial \nabla_{t}\big/\partial\phi = 0$.
%This has the effect of ignoring how changing $\phi$ would have affected earlier
%points of the trajectory when computing the update at time $t$.
This assumption allows us to avoid computing second
derivatives of $f$.

\begin{figure}
\includegraphics[width=1.0\linewidth]{figures/computational-graph}
\caption{Computational graph used for computing the gradient of the optimizer.}
\label{fig:graph}
\end{figure}

% for lengthy
% trajectories, training the optimizer is difficult for two reasons:
% {
% \setlength{\topsep}{0pt}
% \setlength{\itemsep}{0pt}
% \begin{enumerate}

%     \item Training is computationally intensive because we must fully optimize
%     any given $h$ to obtain a single update for the optimizer parameters
%     $\phi$.

%     \item Training is memory intensive because we implement the optimizer with
%     a recurrent neural network and unrolling over a trajectory requires storing
%     the activations for every step \citep[but see][]{maclaurin:2015,chen:2016}.
% \end{enumerate}}
% These are generic problems encountered when training recurrent networks on long
% sequences, both of which are typically addressed by using truncated
% backpropagation through time (BPTT). However, in our setting we see that

Examining the objective in \eqref{eq:meta-optimization} we see that
the gradient is non-zero only for terms where $w_t \neq 0$. If we
use $w_t = \boldsymbol{1}[t = T]$ to match the original problem, then gradients
of trajectory prefixes are zero and only the final optimization step
provides information for training the optimizer. This renders
Backpropagation Through Time (BPTT) inefficient.
We solve this problem by relaxing the objective such that $w_t > 0$ at
intermediate points along the trajectory.  This changes the objective function,
but allows us to train the optimizer on partial trajectories.
For simplicity, in all our experiments we use $w_t = 1$ for every $t$.

% One could imagine
% several ways of setting the $w_t$ to allow this and in this work we choose the
% somewhat extreme option of setting $w_t = 1$ for all $t$. This modified
% objective can be seen to be rewarding improvement per unit of time (over the
% truncated horizon) rather than just the final objective value.

\subsection{Coordinatewise LSTM optimizer}
\label{sec:coordinatewise}

One challenge in applying RNNs in our setting is that we want to be able to
optimize at least tens of thousands of parameters. Optimizing at this scale with a fully
connected RNN is not feasible as it would require a huge hidden state and an
enormous number of parameters. To avoid this difficulty we will use an
optimizer $m$ which operates coordinatewise on the parameters of the objective
function, similar to other common update rules like RMSprop and ADAM. This
\emph{coordinatewise network architecture} allows us to use a very small network that only looks at a single coordinate to
define the optimizer and share optimizer parameters across different parameters
of the optimizee.

Different behavior on each coordinate is achieved by using separate activations
for each objective function parameter. In addition to allowing us to use a
small network for this optimizer, this setup has the nice effect of making the
optimizer invariant to the order of parameters in the network, since the same
update rule is used independently on each coordinate.

\begin{wrapfigure}{r}{0.5\textwidth}
\centering
\includegraphics[width=\linewidth,clip]{figures/lstm-optimizer.pdf}
\caption{One step of an LSTM optimizer. All LSTMs have shared
parameters, but separate hidden states.}
\label{fig:lstm}
\end{wrapfigure}

We implement the update rule for each coordinate using a two-layer Long Short
Term Memory (LSTM) network \citep{hochreiter:1997}, using the now-standard
forget gate architecture. The network takes as input the optimizee gradient for
a single coordinate as well as the previous hidden state and outputs the update
for the corresponding optimizee parameter. We will refer to this architecture,
illustrated in Figure~\ref{fig:lstm}, as an LSTM optimizer.

%As noted above, each LSTM node in this figure shares the same set of parameters
%and learning of the optimizer is performed by unrolling all LSTMs over several
%time steps and computing gradients via backpropagation through time (BPTT).
The use of recurrence allows the LSTM to learn dynamic update rules which
integrate information from the history of gradients, similar to momentum. This
is known to have many desirable properties in convex optimization
\citep[see e.g.][]{nesterov:1983} and in fact many recent learning
procedures---such as ADAM---use momentum in their updates.
%This can be seen by
%examining the flow of forward flow of optimizee derivatives in the
%computational graph.

\paragraph{Preprocessing and postprocessing}

Optimizer inputs and outputs can have very different magnitudes depending on
the class of function being optimized, but neural networks usually work
robustly only for inputs and outputs which are neither very small nor very
large. In practice rescaling inputs and outputs of an LSTM optimizer using
suitable constants (shared across all timesteps and functions $f$) is
sufficient to avoid this problem. In Appendix~\ref{sec:pre} we propose a
different method of preprocessing inputs to the optimizer inputs which is more
robust and gives slightly better performance.


